% Choose onecolumn (default) or twocolumn.
\documentclass[onecolumn]{baysheet}

\title{Programming with clarity}
\author{Your full name}
\date{\today}
\sheetheader{Programming with clarity}

\begin{document}
\maketitle

\section{Pure functions}
A pure function returns the same result for the same input and has no
observable side effects. This sheet uses the same fonts, colors, boxes,
and code listings as the presentation template. Change the class option
to \texttt{twocolumn} to set the body in two columns; the title still spans
the full page width.

\begin{definitionbox}{Pure function}
The result depends only on the function's inputs. Evaluating the function
does not change any state outside the function.
\end{definitionbox}

\subsection{Working with a list}
For the input $[1,2,3,4]$, square each value and then add the results.
The sum is $1+4+9+16=30$. Inline identifiers such as
\texttt{sum\_of\_squares} use Fira Mono.

\begin{examplebox}{Squaring a number}
The function $f(x)=x^2$ maps $4$ to $16$ every time it is evaluated.
\end{examplebox}

\begin{hintbox}{Break it down}
Square each number before adding the results. You can check the result
by calculating the squares separately.
\end{hintbox}

Read more about the language in the
\href{https://docs.python.org/3/tutorial/}{Python tutorial}.

\begin{figure}[htbp]
\centering
\fbox{\rule{0pt}{12mm}\rule{.65\linewidth}{0pt}}
\caption{Replace this placeholder with your own figure.}
\end{figure}

\section{Code and practice}
From the second page onward, the title appears in the running header.
Every page has its page number in the footer.
Listings can use their own line numbers, separate from the prose line
numbers, by passing \texttt{numbers=left}.

\begin{python}[numbers=left]
def sum_of_squares(values):
    return sum(x * x for x in values)

print(sum_of_squares([1, 2, 3, 4]))
\end{python}

\begin{java}
int square(int x) {
    return x * x;
}
\end{java}

\begin{haskell}
square :: Integer -> Integer
square x = x * x
\end{haskell}

\subsection{Try it yourself}
Compute the sum of squares for $[2,3,5]$. Use the same steps as in the
example on the first page, and compare your result with the solution.

\begin{solutionbox}{Worked answer}
The squares are $4$, $9$, and $25$. Adding them gives
$2^2+3^2+5^2=4+9+25=38$.
\end{solutionbox}

Explain why the function is pure and how you would test it with an empty
list. A good explanation distinguishes the input, the returned value,
and any effects on the surrounding program.
\end{document}
